\documentclass[12pt]{article}

\usepackage{amsmath, amssymb}
\usepackage{geometry}
\geometry{margin=1in}

\title{Lecture Scribe -- Lecture 8 (CSE400)}
\author{}
\date{}

\begin{document}

\maketitle

\section*{Topic: Gaussian Random Variable, Standard Forms, and Relation between $\Phi$-function and Q-function}

\section*{1) Definitions and Notations}

\textbf{Definition 1.1 (Gaussian Random Variable)}

A Gaussian random variable is one whose PDF can be written in the general form:

\[
f_X(x)
=
\frac{1}{\sqrt{2\pi\sigma^2}}
\exp
\left(
-
\frac{(x-m)^2}{2\sigma^2}
\right)
\]

\textbf{Notation}

Mean:
\[
m
\]

Variance:
\[
\sigma^2
\]

Standard deviation:
\[
\sigma
\]

CDF:
\[
F_X(x)
\]

Gaussian CDF relation:

\[
F_X(x)
=
\int_{-\infty}^{x}
\frac{1}{\sqrt{2\pi\sigma^2}}
\exp
\left(
-
\frac{(y-m)^2}{2\sigma^2}
\right)
dy
\]

Standard normal CDF ($\Phi$-function):

\[
\Phi(x)
=
\frac{1}{\sqrt{2\pi}}
\int_{-\infty}^{x}
\exp
\left(
-
\frac{t^2}{2}
\right)
dt
\]

Gaussian tail function (Q-function):

\[
Q(x)
=
\frac{1}{\sqrt{2\pi}}
\int_{x}^{\infty}
\exp
\left(
-
\frac{t^2}{2}
\right)
dt
\]

Error function:

\[
\mathrm{erf}(x)
=
\frac{2}{\sqrt{\pi}}
\int_{0}^{x}
\exp(-t^2)
dt
\]

Complementary error function:

\[
\mathrm{erfc}(x)
=
1 - \mathrm{erf}(x)
=
\frac{2}{\sqrt{\pi}}
\int_{x}^{\infty}
\exp(-t^2)
dt
\]

Key identity:

\[
Q(x)
=
1 - \Phi(x)
\]

\section*{2) Assumptions and Preliminaries}

CDF evaluation:

\[
F_X(x)
=
\int_{-\infty}^{x}
\frac{1}{\sqrt{2\pi\sigma^2}}
\exp
\left(
-
\frac{(y-m)^2}{2\sigma^2}
\right)
dy
\]

Standardization:

Let

\[
t =
\frac{y-m}{\sigma}
\]

Then,

\[
F_X(x)
=
\int_{-\infty}^{\frac{x-m}{\sigma}}
\frac{1}{\sqrt{2\pi}}
\exp
\left(
-
\frac{t^2}{2}
\right)
dt
\]

\section*{3) Main Result}

Gaussian CDF relation:

\[
F_X(x)
=
\Phi
\left(
\frac{x-m}{\sigma}
\right)
\]

Tail probability relation:

\[
\Pr(X > x)
=
\int_{\frac{x-m}{\sigma}}^{\infty}
\frac{1}{\sqrt{2\pi}}
\exp
\left(
-
\frac{t^2}{2}
\right)
dt
=
Q
\left(
\frac{x-m}{\sigma}
\right)
\]

Key identities:

\[
Q(x)
=
1 - \Phi(x)
\]

\[
F_X(x)
=
1 -
Q
\left(
\frac{x-m}{\sigma}
\right)
\]

\section*{Proof}

Evaluating the CDF:

\[
F_X(x)
=
\int_{-\infty}^{x}
\frac{1}{\sqrt{2\pi\sigma^2}}
\exp
\left(
-
\frac{(y-m)^2}{2\sigma^2}
\right)
dy
\]

Substitute:

\[
t =
\frac{y-m}{\sigma}
\]

Then,

\[
F_X(x)
=
\int_{-\infty}^{\frac{x-m}{\sigma}}
\frac{1}{\sqrt{2\pi}}
\exp
\left(
-
\frac{t^2}{2}
\right)
dt
=
\Phi
\left(
\frac{x-m}{\sigma}
\right)
\]

Tail probability:

\[
\Pr(X > x)
=
\int_{\frac{x-m}{\sigma}}^{\infty}
\frac{1}{\sqrt{2\pi}}
\exp
\left(
-
\frac{t^2}{2}
\right)
dt
=
Q
\left(
\frac{x-m}{\sigma}
\right)
\]

\section*{Conclusion}

\[
F_X(x)
=
\Phi
\left(
\frac{x-m}{\sigma}
\right)
\]

\[
\Pr(X > x)
=
Q
\left(
\frac{x-m}{\sigma}
\right)
\]

\[
Q(x)
=
1 - \Phi(x)
\]

\section*{4) Working Examples}

\textbf{Example 1}

Given PDF:

\[
f_X(x)
=
\frac{1}{\sqrt{8\pi}}
\exp
\left(
-
\frac{(x+3)^2}{8}
\right)
\]

Comparing with Gaussian PDF:

\[
m = -3
\]

\[
\sigma^2 = 4
\]

\[
\sigma = 2
\]

Gaussian relations:

\[
F_X(x)
=
\Phi
\left(
\frac{x-m}{\sigma}
\right)
\]

\[
\Pr(X>x)
=
Q
\left(
\frac{x-m}{\sigma}
\right)
\]

\textbf{(a) Find $\Pr(X \le 0)$}

\[
\Pr(X \le 0)
=
F_X(0)
\]

\[
=
\Phi
\left(
\frac{0-(-3)}{2}
\right)
\]

\[
=
\Phi
\left(
\frac{3}{2}
\right)
\]

\[
=
1 -
Q
\left(
\frac{3}{2}
\right)
\]

\textbf{(b) Find $\Pr(X > 4)$}

\[
\Pr(X > 4)
=
Q
\left(
\frac{4-(-3)}{2}
\right)
\]

\[
=
Q
\left(
\frac{7}{2}
\right)
\]

\end{document}
